\section{Proofs for Section~\ref{sec:comparative}}
\label{app:comparative}

\begin{proof}[Proof of Proposition~\ref{prop:fewer-rights}]
The upper bound holds because $\pi_{\cK}$ is admissible in $\mathcal P'$. For
the lower bound, induct on $L(S)$ over corresponding states, which share
the own states and sunk cost $K$ and have $q\le q'$. If $\pi_{\cK}$ stops at $\bar S$ then $\max_iG_i(s_i)\le q'$, and since $Z_i\le G_i$ pathwise,
Theorem~\ref{thm:index} gives
$V^\star(S)+K=\E[\max\{q,\max_iZ_i\}]\le\max\{q,\max_iG_i\}\le q'$, so
$V^\star(S)\le q'-K=V_{\pi_{\cK}}(\bar S)$. Otherwise $\pi_{\cK}$ advances an arm
$i$ with $G_i(s_i)=\max_jG_j(s_j)>q'\ge q$, which is an optimal action in
$\mathcal P$ at $S$ by Theorem~\ref{thm:index}. Both problems pay $c(s_i)$ and
reveal the same reading $Y$; the successors correspond, with incumbents
$\max\{q,\Yc\}\le\max\{q',\Yc'\}$ since $\Yc\le\Yc'$. Hence
$V^\star(S)=\E[V^\star(S_1)]\le\E[V_{\pi_{\cK}}(\bar S_1)]
=V_{\pi_{\cK}}(\bar S)$ by the induction hypothesis applied under the integral.
The successor values already include the advance cost through
$K(S_1)=K(\bar S_1)=K+c(s_i)$.
Appendix~\ref{app:comparative-results} exhibits cells attaining each bound.
\end{proof}

\begin{proof}[Proof of Proposition~\ref{prop:resumption}]
(i) The index policy is no-return whenever it never resumes, which gives the
equality case. Let $\pi$ be any no-return policy; conditional on its auxiliary randomness it is a deterministic no-return policy, so we may take it deterministic. Let $\sigma$ be the first
step at which $\pi$'s action differs from $\pi_{\rm IX}$'s at the state
reached, with $\sigma=\infty$ if they never differ. Up to $\sigma$ both
policies take the same actions and observe the same readings, so the law of
$(\widehat S_0,\ldots,\widehat S_\sigma)$ under $\pi$ is the law of the same stopped segment of
the $\pi_{\rm IX}$ trajectory, including the last-arm and left-arm bookkeeping
from the fixed root $S_0$. Since $\pi$ cannot resume, $\sigma$ is at most
the first time $\pi_{\rm IX}$ resumes. Theorem~\ref{thm:loss} writes
$V^\star(S_0)-V_\pi(S_0)$ as the expected sum of nonnegative losses along $\pi$,
one of which, on $\{\sigma<\infty\}$, is the loss of $\pi$'s action at
$S_\sigma$; that action differs from $\pi_{\rm IX}(S_\sigma)$, so its loss is at
least $\delta(S_\sigma)$. Hence
$V^\star(S_0)-V_\pi(S_0)\ge\E_{\pi_{\rm IX}}[\delta(S_\sigma)\mathbf 1\{\sigma<\infty\}]$.
The right side is the cost of a stopping time on the $\pi_{\rm IX}$ trajectory
that must stop no later than the first resumption and may run forever only on
branches without resumption. The infimum of this cost over such stopping times
satisfies the finite-horizon recursion defining $\widehat\rho$: it is zero where no
resumption can occur, equals $\delta$ at a resumption state, and otherwise is
the smaller of stopping now and continuing. Thus
$V^\star(S_0)-V_\pi(S_0)\ge\widehat\rho(S_0,0,\varnothing)=\rho(S_0)$
for every no-return $\pi$.

(ii) At a state $S$ where the index-maximal arm $i^\star$ is unique with
$G_{i^\star}>q$, every other action has positive loss. By
Lemma~\ref{lem:nested}, $\E[(W_{i^\star}-G_{i^\star})^+]=c(s_{i^\star})>0$,
so $W_{i^\star}>G_{i^\star}$ and hence $Z_{i^\star}=G_{i^\star}$ with positive
probability. Then $\Delta_{\rm stop}(S)=\E[(\max_jZ_j-q)^+]>0$. For an arm
$j\ne i^\star$ with finite index, $O_j\ge Z_{i^\star}=G_{i^\star}>G_j$ with
positive probability, so $\Prob(O_j>G_j)>0$ and
$\Delta_j(S)>0$ by the paragraph following Theorem~\ref{thm:loss}; if $G_j=-\infty$
then $\Delta_j(S)=c(s_j)>0$. Hence $\delta(S)>0$ at every state in the
hypothesis. The augmented recursion then gives $\widehat\rho(\widehat S)>0$
if and only if resumption from that augmented history has positive probability:
at a resumption state $\widehat\rho=\delta>0$, on branches without resumption
$\widehat\rho=0$, and elsewhere
$\widehat\rho=\min\{\delta,\E[\widehat\rho(\widehat S_1)]\}$ is positive exactly when
$\E[\widehat\rho(\widehat S_1)]>0$. Finite ladders make the induction terminate;
evaluating at $(S_0,0,\varnothing)$ proves the assertion for $\rho(S_0)$.

(iii) Let every reading be a deterministic function of the own state. Then
$W(s)$ is a number, the reservation equation $(W(s)-g)^+=c(s)$ has the unique
solution $G(s)=W(s)-c(s)$ when a collectable stage remains, $Z(s)=G(s)$, and
$W(s)=\max\{\Yc(s),G(s')\}$ with $s'$ the successor. Consequently an advance
from $s$ to $s'$ either has $\Yc(s)\ge G(s')$, in which case the new incumbent
is at least $\Yc(s)=G(s)+c(s)>G(s)$, or has $\Yc(s)<G(s')$, in which case
$G(s')=G(s)+c(s)>G(s)$; the first case is forced when $s'$ is exhausted.
Suppose $\pi_{\rm IX}$ advances an arm $i\ne j$ at a state where arm $j$, which it has advanced from $S_0$, is not exhausted; then $G_i(s_i)\ge G_j(s_j)$ and $G_i(s_i)>q$.
We claim that at every later state either $q>G_j(s_j)$ or some unexhausted
arm $k$ has $G_k(s_k)>G_j(s_j)$. After the advance of $i$, the first case
above gives $q>G_i(s_i)\ge G_j(s_j)$ and the second gives
$G_i(s_i')>G_i(s_i)\ge G_j(s_j)$. If the claim holds with $q>G_j(s_j)$ it holds
forever, as the incumbent never decreases and $s_j$ does not change while $j$
is paused. If it holds through an arm $k$, the next advance, if any, is of an
arm $m$ with $G_m(s_m)\ge G_k(s_k)>G_j(s_j)$, so $m\ne j$, and the two cases
above restore the claim with $q>G_m(s_m)>G_j(s_j)$ or with $k=m$. In neither
case is $j$ ever index-maximal with index above the incumbent, so it is never
resumed, whatever the tie order.
\end{proof}

\begin{example}[Resumption without learning future rewards]\label{ex:resumption-reward}
Let $q_0=0$, make every stage collectable, and charge $1/10$ per advance.
Arm $A$ first pays zero or one with equal probability, then pays $2/5$
deterministically; independent arm $B$ pays zero or $3/5$ with equal probability
in one stage. The indices are $G_{A,1}=4/5$, $G_B=2/5$, and $G_{A,2}=3/10$.
The index policy starts $A$, switches to $B$ after a zero, and resumes $A$ after
a second zero. Its value is $23/40$, whereas the no-return optimum is $11/20$,
so resumption is worth $1/40$. No reading changes any remaining reward law:
$A$'s index drops at its known stage transition, while the realized collectable
reward determines whether to continue. The calculation below verifies both values.
\end{example}

\paragraph{Value calculations.}
The index policy stops after a first reward of one; after a zero from $A$ it
tries $B$ and completes $A$ only if $B$ also returns zero. Its net utility is
\[
 \tfrac12(1-\tfrac1{10})
 +\tfrac14(\tfrac35-\tfrac2{10})
 +\tfrac14(\tfrac25-\tfrac3{10})=\tfrac{23}{40}.
\]
For a no-return policy starting $A$, a zero leaves the choice between completing
$A$ for additional value $3/10$ and abandoning it for $B$, whose value is $1/5$.
Completing $A$ raises the incumbent to $2/5$, at which inspecting $B$ ties stopping.
Thus starting $A$ has no-return value
$-1/10+(1/2)\cdot1+(1/2)\cdot(3/10)=11/20$.
Starting $B$ and then acting optimally on $A$ gives $21/40$, and stopping gives zero;
these exhaust the first actions and prove the no-return optimum.
